\subsection{Why an infinite projective arithmetic quotient does not remove ramification}

\begin{result}{A.32}{Proved}{A projective quotient retaining the cyclotomic character kills off-diagonal tame inertia}
Let \(K\) be a number field, \(\ell\) a prime, and
\(\psi:G_K\twoheadrightarrow\Gamma\) a continuous quotient.
Assume that \(\Gamma\) has an open pro-\(\ell\) subgroup,
that \(\operatorname{cd}_\ell\Gamma\leq1\), and that the
\(\ell\)-adic cyclotomic character factors continuously through
\(\psi\). Then, at each finite place \(v\nmid\ell\), the
\(\ell\)-primary tame inertia has trivial image in \(\Gamma\)
after passage to any finite extension making the local image
pro-\(\ell\).

In particular, suppose that the actual tangential action on
\(F_\ell\), and a continuous section cocycle
\(G_{\mathbb Q}\to F_\ell\), factor through such a quotient of
\(G_{\mathbb Q}\). Its abelian pair is unramified away from
\(\ell\). If that pair is the \(\ell\)-primary component of an
ordinary completed pair, the latter has finite support and is a
rational unit-equation pair. This holds without a finite-ramification
assumption on \(\Gamma\).
\end{result}
\begin{proof}
Fix \(v\nmid\ell\). Replace \(K_v\) by a finite extension
so that its image lies in an open pro-\(\ell\) subgroup of
\(\Gamma\), and so that its residue cardinality \(Q\) satisfies
\(Q\equiv1\pmod\ell\), or \(Q\equiv1\pmod4\) if
\(\ell=2\). This replacement is allowed because the image of
inertia under a finite extension is open in its previous image.
The wild inertia, being pro-residue-characteristic, has trivial
image in a pro-\(\ell\) group. The remaining inertia image
\(T\) is therefore a quotient of \(\mathbb Z_\ell(1)\).

A profinite group of finite \(\ell\)-cohomological dimension
has no nontrivial finite \(\ell\)-subgroup. Indeed such a subgroup
has nonzero \(\mathbb F_\ell\)-cohomology in arbitrarily high
degrees, whereas coinduction and Shapiro's identity give
\(\operatorname{cd}_\ell H\leq\operatorname{cd}_\ell\Gamma\)
for every closed subgroup \(H\). Thus, if \(T\ne1\), then
\(T\simeq\mathbb Z_\ell\).

Let \(h\) be the image of arithmetic Frobenius. Its closed cyclic
subgroup is pro-\(\ell\), and its cyclotomic image is generated
by \(Q\). This element has infinite order in
\(1+\ell\mathbb Z_\ell\), respectively
\(1+4\mathbb Z_2\): a positive rational integer \(Q>1\)
cannot be a root of unity in \(\mathbb Q_\ell\).
Consequently \(\overline{\langle h\rangle}\simeq\mathbb Z_\ell\)
and its cyclotomic character is injective. The cyclotomic character
is trivial on \(T\), so these two closed subgroups intersect
trivially. The tame relation gives
\[
 h\tau h^{-1}=\tau^Q
 \quad(\tau\text{ a topological generator of }T).
\]
Their product is compact and is a closed subgroup of \(\Gamma\)
isomorphic to
\[
       H=\mathbb Z_\ell\rtimes_Q\mathbb Z_\ell.
\]
Both cyclic groups have cohomological dimension one. The
Hochschild--Serre spectral sequence with trivial coefficients
\(\mathbb F_\ell\) has, in total degree two, just
\[
 H^1\bigl(\mathbb Z_\ell,
              H^1(\mathbb Z_\ell,\mathbb F_\ell)\bigr)
       =\mathbb F_\ell,
\]
since \(Q\equiv1\pmod\ell\). There are no possible incoming
or outgoing differentials, so \(H^2(H,\mathbb F_\ell)\ne0\).
This contradicts \(\operatorname{cd}_\ell\Gamma\leq1\).
We conclude that \(T=1\).

For the section consequence, take \(q\ne\ell\). The actual
tangential action is unramified at \(q\), as proved in A.31.
The restriction of its cocycle to inertia is a continuous
homomorphism into the torsion-free pro-\(\ell\) group
\(F_\ell\). After the finite local extension just used, it
vanishes: it factors through \(\psi\), whose pro-\(\ell\)
inertia image is trivial. Before that extension its image is finite,
because the subgroup on which it vanishes has finite index. A finite
subgroup of the free pro-\(\ell\) group is trivial, so the
original inertia cocycle also vanishes. Apply the abelianization
map and then A.3 and E.101. Nonzero ordinary integer valuations
cannot disappear in one \(\ell\)-adic component. Thus the
ordinary pair has support contained in \(\{\ell\}\), in
particular finite support, and is rational.
\end{proof}

\paragraph{Scope of the failed construction.}
The full lifts in P.4 use a virtually pro-\(\ell\) arithmetic
quotient of \(\ell\)-cohomological dimension at most one.
A.32 shows that dropping only its finite-ramification restriction
cannot turn that mechanism into an ordinary ghost construction.
This does not rule out a quotient of cohomological dimension two,
or a non-projective embedding-problem argument. In particular, the
free relative inertia kernel in arithmetic Riemann existence is not
the projective quotient considered here. The quotient action on
that kernel must still satisfy the Frobenius--inertia relations;
its freeness alone does not solve the matching equations.
The precise source statements are
\citet[Proposition~8.1 and Corollary~8.2, pp.~467--468]{RootSchmidt2007}
and \citet[Theorem~1.1 and Remarks~1.2, author p.~2]{RootSchmidt2010};
the latter explicitly excludes coefficient two. Their hypotheses and
the failed arithmetic-extension step are recorded in
\path{notes/main_gate_search_06.md}.

The current family anabelian theorem also does not provide a section
lifting statement:
\citet[Theorem~A, p.~2]{RootHolzschuhSchmidtStix2026} requires
a map open on the full profinite fundamental group. A section has
infinite-index image, since the geometric kernel is infinite, so it
fails that hypothesis. A map to the pronilpotent quotient has a
different target as well. We use only this explicit scope check,
not the general theorem as an input to a section proof.

